\section*{الفصل \ref{ch:g2:balance-equilibrium} --- الاتزان والتوازن}

\begin{solution}{exo:g2:balance-equilibrium:1}
الأرض تسحب المصباح إلى الأسفل؛ والمكتب يدفعه إلى الأعلى بالشدّة
نفسها بالضبط. فتتعادل القوّتان --- \omterm{def:g2:balance-equilibrium:equilibrium}{توازن} --- فيبقى المصباح
مكانه.
\end{solution}

\begin{solution}{exo:g2:balance-equilibrium:2}
حين يسحب الفريقان بالشدّة نفسها في جهتين متعاكستين: فتلغي
السحبتان إحداهما الأخرى. وينكسر التعادل لحظة يسحب أحد الفريقين
بشدّة أكبر من الآخر.
\end{solution}

\begin{solution}{exo:g2:balance-equilibrium:3}
تستوي أفقيًّا --- وزنان متساويان على مسافتين متساويتين. وحين
ينزلق أحد التوأمين مبتعدًا ينزل جانبه: فالمسافة تحسب أيضًا.
\end{solution}

\begin{solution}{exo:g2:balance-equilibrium:4}
قريبًا من المحور، ولينا بعيدة في جانبها. والقاعدة: الراكب الأثقل
قرب الوسط يستطيع أن يتعادل مع راكب أخفّ بعيد --- فالوزن والمسافة
كلاهما يحسب.
\end{solution}

\begin{solution}{exo:g2:balance-equilibrium:5}
\omterm{def:g2:balance-equilibrium:point}{نقطة توازن} المسطرة في وسطها؛ \omterm{def:g2:balance-equilibrium:point}{ونقطة توازن} المطرقة قرب رأسها
الثقيل. وللتحقّق: ضع كلًّا منهما على إصبع واحدة وجِد الموضع
الذي تستقرّ فيه مستوية.
\end{solution}

\begin{solution}{exo:g2:balance-equilibrium:6}
تلتقيان قرب الفرشاة --- أي الطرف الثقيل. \omterm{def:g2:balance-equilibrium:point}{فنقطة التوازن} تقع نحو
الجهة الثقيلة من الجسم.
\end{solution}

\begin{solution}{exo:g2:balance-equilibrium:7}
الهرم العريض: إذ يستطيع أن يميل ميلًا كبيرًا قبل أن يخرج وزنه
عن قاعدته العريضة. أما البرج الضيّق فينقلب عند أدنى ميل.
\end{solution}

\begin{solution}{exo:g2:balance-equilibrium:8}
العصا تُبطئ كل ترنّح، فتمنح الماشية وقتًا لتعود فوق الحبل. وعلى
حافة الرصيف تمدّ ذراعيك للسبب نفسه: ترنّح أبطأ، وتدارك أسهل.
\end{solution}

\begin{solution}{exo:g2:balance-equilibrium:9}
لتقف على القدم اليمنى وحدها لا بدّ أن يميل جسمك إلى اليمين، حتى
يقع وزنك فوق تلك القدم --- مثل البرج الذي يبقي وزنه فوق قاعدته.
والجدار في الموضع الذي تحتاج إلى الميل إليه بالضبط: فهو يمنع
الحركة، ويبقى وزنك معلّقًا فوق موضع القدم المرفوعة الخالي،
فتميل وتسقط.
\end{solution}

\begin{solution}{exo:g2:balance-equilibrium:10}
الرافعة أرجوحة متأرجحة عملاقة، وبرجها هو المحور. فالراكب الأول:
الحمل، بعيدًا على الذراع الأمامية الطويلة. والراكب الثاني: الكتل
الخرسانية --- وهي أثقل بكثير، فيجوز أن تجلس قريبًا، على الذراع
الخلفية القصيرة. فالثقيل القريب يتعادل مع الخفيف البعيد، وهي
قاعدة الأرجوحة بعينها؛ والذراع الخلفية قصيرة لأن راكبها هو
الوزن الثقيل.
\end{solution}
